El proyecto deja abiertas cuatro líneas de evolución, ordenadas por
prioridad y respaldadas por la arquitectura actual del sistema.

\subsection{Endurecimiento para despliegue en producción}

El \cref{rnf:deployment} se satisface con el fichero
\texttt{docker-compose.yml} actual, pero varios aspectos deben reforzarse
antes de un despliegue en producción. La sustitución del servidor de
desarrollo de Django (\texttt{runserver}) por un servidor \acs{wsgi}
dedicado es la primera medida, y resolvería el cuello de
botella detectado en \cref{SEC:TEST-LOAD}. La
sustitución del \textit{watcher} sobre volumen Docker por un cliente
\acs{sftp} real conectado a cabinas de escaneo completaría la preparación para producción.

\subsection{Integraciones externas}

La arquitectura de \dfnpl{plugin} descrita en \cref{SEC:DES-PLUGINS}
facilita tres integraciones previstas pero no implementadas: la
activación de \acs{sso} institucional (\acs{saml}~2.0 y \acs{oidc}) sobre
el \dfn{plugin} de autenticación ya preparado; la
conexión con \acsp{lms} mediante \acs{lti}~1.3, que abriría el sistema a
Moodle y otras plataformas del entorno universitario.

\subsection{Funcionalidad ampliada}

Varias líneas de mejora funcional son viables sobre la base actual. Habilitar
la ejecución asíncrona de las importaciones masivas mediante tareas Celery.
Generar estadísticas comparativas entre convocatorias del mismo examen. O ampliar
las opciones de annotación (pngs, gifs, etc.) completarían el sistema.

\subsection{Refinamientos operativos y de experiencia de usuario}

El último bloque agrupa mejoras de operación y usabilidad: extender todavía
más los parámetros configurables por \dfn{tenant} y por ejemplo, la introducción
de notificaciones en tiempo real mediante WebSockets sobre
\acs{asgi} para eliminar la dependencia del \textit{polling} actual.
